% GENERATED FILE. Edit canonical lessons, then run npm run generate:books.
% canonical-source-hash: fnv1a64:bba5e3e60cb33d32
% canonical-lessons: SA-C381-samanyatah, SA-C381-visesatah, SA-C381-sarvatha, SA-C381-bahusah, SA-C381-sanksipati

\chapter{Usually and Especially}
\label{ch:sa-usually-and-especially}
\hlchaptermodality{381}

\begin{chapteropening}
\textbf{By the end of this chapter:} \emph{I can say what usually happens and what happens especially.}

Complete the last lesson of chapter 381: I can say what usually happens and what happens especially.
\end{chapteropening}

\section[\texorpdfstring{\sk{सामान्यतः} (sāmānyataḥ)}{सामान्यतः (sāmānyataḥ)}]{\sk{सामान्यतः} (sāmānyataḥ) — generally; as a rule}

\label{lesson:SA-C381-samanyatah}



\begin{quote}
Before the new one: say the Sanskrit for as much as one likes, then the Sanskrit for expands.
\end{quote}

\subsection*{You'll want to know: \sk{सामान्यतः}}
\textbf{\sk{सामान्यतः}} — \emph{sāmānyataḥ} — \textquotedblleft{}generally; as a rule\textquotedblright{}.

\begin{quote}
\sk{सामान्यतः}, \sk{यदि} \sk{अहं} \sk{श्रान्तः}, \sk{तर्हि} \sk{शीघ्रं} \sk{स्वपिमि।}
\end{quote}

\begin{quote}
\emph{sāmānyataḥ, yadi ahaṃ śrāntaḥ, tarhi śīghraṃ svapimi.}
\end{quote}

\begin{quote}
Generally, if I am tired, I fall asleep quickly.
\end{quote}

\begin{cousinweb}
\textbf{\sk{सामान्यतः}} (\emph{sāmānyataḥ}) is built on \emph{sāmānya}, \textquotedblleft{}common, general\textquotedblright{}, with the ending \emph{-tas}; MW gives \emph{sāmānya-tas} \textquotedblleft{}in general, generally\textquotedblright{} (MW).
\end{cousinweb}

\begin{grammarlens}[title={What you've built}]
One more word for saying what happens if.
\end{grammarlens}

\subsection*{Your turn}
\begin{itemize}
\raggedright
  \item \emph{Say it:} \emph{sāmānyataḥ}
  \item \emph{Say it:} \emph{sāmānyataḥ}, once more
  \item \emph{Say it:} \emph{sāmānyataḥ, yadi ahaṃ śrāntaḥ, tarhi śīghraṃ svapimi.} — Generally, if I am tired, I fall asleep quickly.
  \item \emph{Recall:} say the Sanskrit for as much as one likes, then the Sanskrit for expands, then say \emph{sāmānyataḥ} again
\end{itemize}

\begin{tcolorbox}[breakable,colback=teal!4,colframe=teal!35!black,title={Before you move on}]
What does \sk{सामान्यतः} mean? (\textquotedblleft{}Generally; as a rule\textquotedblright{}.) Say it once more.
\end{tcolorbox}
\section[\texorpdfstring{\sk{विशेषतः} (viśeṣataḥ)}{विशेषतः (viśeṣataḥ)}]{\sk{विशेषतः} (viśeṣataḥ) — especially}

\label{lesson:SA-C381-visesatah}



\begin{quote}
Before the new one: say the Sanskrit for expands, then the Sanskrit for generally.
\end{quote}

\subsection*{You'll want to know: \sk{विशेषतः}}
\textbf{\sk{विशेषतः}} — \emph{viśeṣataḥ} — \textquotedblleft{}especially\textquotedblright{}.

\begin{quote}
\sk{यदि} \sk{मार्गः} \sk{आर्द्रः}, \sk{तर्हि} \sk{विशेषतः} \sk{मन्दं} \sk{चलतु।}
\end{quote}

\begin{quote}
\emph{yadi mārgaḥ ārdraḥ, tarhi viśeṣataḥ mandaṃ calatu.}
\end{quote}

\begin{quote}
If the road is wet, walk especially slowly.
\end{quote}

\begin{cousinweb}
\textbf{\sk{विशेषतः}} (\emph{viśeṣataḥ}) is built on \emph{viśeṣa}, \textquotedblleft{}a distinction\textquotedblright{}, with the ending \emph{-tas}; MW gives \emph{viśeṣa-tas} \textquotedblleft{}especially, particularly, above all\textquotedblright{} (MW).
\end{cousinweb}

\begin{grammarlens}[title={What you've built}]
One more word for saying what happens if.
\end{grammarlens}

\subsection*{Your turn}
\begin{itemize}
\raggedright
  \item \emph{Say it:} \emph{viśeṣataḥ}
  \item \emph{Say it:} \emph{viśeṣataḥ}, once more
  \item \emph{Say it:} \emph{yadi mārgaḥ ārdraḥ, tarhi viśeṣataḥ mandaṃ calatu.} — If the road is wet, walk especially slowly.
  \item \emph{Recall:} say the Sanskrit for expands, then the Sanskrit for generally, then say \emph{viśeṣataḥ} again
\end{itemize}

\begin{tcolorbox}[breakable,colback=teal!4,colframe=teal!35!black,title={Before you move on}]
What does \sk{विशेषतः} mean? (\textquotedblleft{}Especially\textquotedblright{}.) Say it once more.
\end{tcolorbox}
\section[\texorpdfstring{\sk{सर्वथा} (sarvathā)}{सर्वथा (sarvathā)}]{\sk{सर्वथा} (sarvathā) — in every way; altogether}

\label{lesson:SA-C381-sarvatha}



\begin{quote}
Before the new one: say the Sanskrit for generally, then the Sanskrit for especially.
\end{quote}

\subsection*{You'll want to know: \sk{सर्वथा}}
\textbf{\sk{सर्वथा}} — \emph{sarvathā} — \textquotedblleft{}in every way; altogether\textquotedblright{}.

\begin{quote}
\sk{यदि} \sk{भवान्} \sk{आगच्छति}, \sk{तर्हि} \sk{अहं} \sk{सर्वथा} \sk{साहाय्यं} \sk{करिष्यामि।}
\end{quote}

\begin{quote}
\emph{yadi bhavān āgacchati, tarhi ahaṃ sarvathā sāhāyyaṃ kariṣyāmi.}
\end{quote}

\begin{quote}
If you come, I will help in every way.
\end{quote}

\begin{cousinweb}
\textbf{\sk{सर्वथा}} (\emph{sarvathā}) is built on \emph{sarva}, \textquotedblleft{}all\textquotedblright{}, with the ending \emph{-thā}, \textquotedblleft{}in the manner of\textquotedblright{}; MW gives \emph{sarvathā} \textquotedblleft{}in every way, by all means, altogether\textquotedblright{} (MW).
\end{cousinweb}

\begin{grammarlens}[title={What you've built}]
One more word for saying what happens if.
\end{grammarlens}

\subsection*{Your turn}
\begin{itemize}
\raggedright
  \item \emph{Say it:} \emph{sarvathā}
  \item \emph{Say it:} \emph{sarvathā}, once more
  \item \emph{Say it:} \emph{yadi bhavān āgacchati, tarhi ahaṃ sarvathā sāhāyyaṃ kariṣyāmi.} — If you come, I will help in every way.
  \item \emph{Recall:} say the Sanskrit for generally, then the Sanskrit for especially, then say \emph{sarvathā} again
\end{itemize}

\begin{tcolorbox}[breakable,colback=teal!4,colframe=teal!35!black,title={Before you move on}]
What does \sk{सर्वथा} mean? (\textquotedblleft{}In every way; altogether\textquotedblright{}.) Say it once more.
\end{tcolorbox}
\section[\texorpdfstring{\sk{बहुशः} (bahuśaḥ)}{बहुशः (bahuśaḥ)}]{\sk{बहुशः} (bahuśaḥ) — often; repeatedly}

\label{lesson:SA-C381-bahusah}



\begin{quote}
Before the new one: say the Sanskrit for especially, then the Sanskrit for in every way.
\end{quote}

\subsection*{You'll want to know: \sk{बहुशः}}
\textbf{\sk{बहुशः}} — \emph{bahuśaḥ} — \textquotedblleft{}often; repeatedly\textquotedblright{}.

\begin{quote}
\sk{यदि} \sk{भवान्} \sk{बहुशः} \sk{अभ्यस्यति}, \sk{तर्हि} \sk{कुशलः} \sk{भविष्यति।}
\end{quote}

\begin{quote}
\emph{yadi bhavān bahuśaḥ abhyasyati, tarhi kuśalaḥ bhaviṣyati.}
\end{quote}

\begin{quote}
If you practise often, you will become skilled.
\end{quote}

\begin{cousinweb}
\textbf{\sk{बहुशः}} (\emph{bahuśaḥ}) is built on \emph{bahu}, \textquotedblleft{}much, many\textquotedblright{}, with the ending \emph{-śas}; MW gives \emph{bahu-śas} \textquotedblleft{}manifoldly, repeatedly, much, often\textquotedblright{} (MW).
\end{cousinweb}

\begin{grammarlens}[title={What you've built}]
One more word for saying what happens if.
\end{grammarlens}

\subsection*{Your turn}
\begin{itemize}
\raggedright
  \item \emph{Say it:} \emph{bahuśaḥ}
  \item \emph{Say it:} \emph{bahuśaḥ}, once more
  \item \emph{Say it:} \emph{yadi bhavān bahuśaḥ abhyasyati, tarhi kuśalaḥ bhaviṣyati.} — If you practise often, you will become skilled.
  \item \emph{Recall:} say the Sanskrit for especially, then the Sanskrit for in every way, then say \emph{bahuśaḥ} again
\end{itemize}

\begin{tcolorbox}[breakable,colback=teal!4,colframe=teal!35!black,title={Before you move on}]
What does \sk{बहुशः} mean? (\textquotedblleft{}Often; repeatedly\textquotedblright{}.) Say it once more.
\end{tcolorbox}
\section[\texorpdfstring{\sk{संक्षिपति} (saṅkṣipati)}{संक्षिपति (saṅkṣipati)}]{\sk{संक्षिपति} (saṅkṣipati) — sums up; shortens}

\label{lesson:SA-C381-sanksipati}



\begin{quote}
Before the new one: say the Sanskrit for in every way, then the Sanskrit for repeatedly.
\end{quote}

\subsection*{You'll want to know: \sk{संक्षिपति}}
\textbf{\sk{संक्षिपति}} — \emph{saṅkṣipati} — \textquotedblleft{}sums up; shortens\textquotedblright{}.

\begin{quote}
\sk{यदि} \sk{समयः} \sk{अल्पः}, \sk{तर्हि} \sk{वक्ता} \sk{भाषणं} \sk{संक्षिपति।}
\end{quote}

\begin{quote}
\emph{yadi samayaḥ alpaḥ, tarhi vaktā bhāṣaṇaṃ saṅkṣipati.}
\end{quote}

\begin{quote}
If time is short, the speaker shortens the speech.
\end{quote}

Past (he …): \sk{सः} \sk{भाषणं} \sk{संक्षिप्तवान्}, \emph{saḥ bhāṣaṇaṃ saṅkṣiptavān}.

\begin{cousinweb}
\textbf{\sk{संक्षिपति}} (\emph{saṅkṣipati}) is built on \emph{sam}, \textquotedblleft{}together\textquotedblright{}, and the root √kṣip, \textquotedblleft{}to throw\textquotedblright{}; MW gives \emph{saṃ-√kṣip} \textquotedblleft{}to throw or heap together\textquotedblright{} and \textquotedblleft{}to concentrate\textquotedblright{}, and \textquotedblleft{}to shorten\textquotedblright{} is how modern Sanskrit uses it (MW).
\end{cousinweb}

\begin{grammarlens}[title={What you've built}]
One more word for saying what happens if.
\end{grammarlens}

\subsection*{Your turn}
\begin{itemize}
\raggedright
  \item \emph{Say it:} \emph{saṅkṣipati}
  \item \emph{Say it:} \emph{saṅkṣipati}, once more
  \item \emph{Say it:} \emph{yadi samayaḥ alpaḥ, tarhi vaktā bhāṣaṇaṃ saṅkṣipati.} — If time is short, the speaker shortens the speech.
  \item \emph{Recall:} say the Sanskrit for in every way, then the Sanskrit for repeatedly, then say \emph{saṅkṣipati} again
\end{itemize}

\begin{tcolorbox}[breakable,colback=teal!4,colframe=teal!35!black,title={Before you move on}]
What does \sk{संक्षिपति} mean? (\textquotedblleft{}Sums up; shortens\textquotedblright{}.) Say it once more.
\end{tcolorbox}
